% Cote 119: offprints, reports, courses, CV and letters (typed and printed part): archive cote 119, site page 289 (PDF page 290).
% typescript. Draft compared: none.
% AI-transcribed by Claude Opus 5.5 (claude-opus-5-5, adaptive thinking, effort high), Message Batches API, batch r119_full (msgbatch_01DwZNm5Y7YsKBWyAN4sFs8w), 2026-11-13; not yet checked by eye.
\scanpage{119/289}
\mspage{- 2 -}

§2. Le cas orthogonal

C'est celui qui est le mieux compris, grâce à son interprétation en termes
de classes de formes quadratiques :

Soit $q$ une forme quadratique non dégénérée de rang $n \geqslant 1$ sur $k$,
et soit $O(q)$ le \uline{groupe orthogonal} de $q$, vu comme groupe algébrique
sur $k$. Si $x$ est un élément de $H^1(k, O(q))$, on peut \uline{tordre} $q$ par $x$ et
l'on obtient une autre forme quadratique $q_x$ de même rang $n$ que $q$. L'appli-%
cation $x \longmapsto (q_x)$ définit une \uline{bijection} de $H^1(k, O(q))$ sur l'ensemble des
\uline{classes de formes quadratiques non dégénérées de rang} $n$ sur $k$.

On a un résultat analogue pour la composante neutre $SO(q)$ de $O(q)$, à
condition de se borner aux formes quadratiques ayant même discriminant que $q$.

Ainsi, tout \uline{invariant} des classes de formes quadratiques peut être
interprété comme une fonction sur l'ensemble de cohomologie $H^1(k, O(q))$,
ou sur l'ensemble $H^1(k, SO(q))$.

2.1. Exemples d'invariants : les classes de Stiefel-Whitney

Ecrivons $q$ comme somme directe orthogonale \xout{\illegible{1}} de formes de
rang $1$ :
\[ q = \langle a_1\rangle \oplus \langle a_2\rangle \oplus \ldots \oplus \langle a_n\rangle = \langle a_1, a_2, \ldots, a_n\rangle , \text{ avec } a_i \in k^* . \]
Si $m$ est un entier $\geqslant 0$ , on définit un élément $w_m(q)$ de $H^m(k, Z/2Z)$
par la formule
\[ \text{(2.1.1)} \qquad w_m(q) = \sum_{i_1 < \ldots < i_m} (a_{i_1})\ldots(a_{i_m}) . \]
(On a noté $(a)$ l'élément de $H^1(k, Z/2Z)$ défini par $a \in k^*$ ; le produit
$(a_{i_1})\ldots(a_{i_m})$ est un cup-produit dans l'algèbre de cohomologie $H^*(k,Z/2Z)$.)

On montre (A.Delzant, \uline{C}.\uline{R}.\uline{Acad}.\uline{Sci}.\uline{Paris}, 255, 1962) que $w_m(q)$ ne dépend
que de la classe d'isomorphisme de $q$ (et pas de la décomposition choisie) ;
cela provient du fait bien connu que les relations entre formes quadratiques
`` résultent des relations en rang $\leqslant 2$ '' .

On dit que $w_m(q)$ est la m-ième \uline{classe de Stiefel-Whitney} de $q$.
\uline{Remarques}. 1) Les classes $w_1(q)$ et $w_2(q)$ ont des interprétations standard:
discriminant, invariant de Hasse-Witt. Les $w_m(q)$ , $m \geqslant \hand{3}$, sont moins intéres-%
santes; il y a avantage à les remplacer (dans la mesure du possible)
par les invariants de la théorie de Milnor, cf. n° 2.3 ci-après.

2) La même méthode conduit à d'autres invariants. Ainsi, si $n$ est

\mspage{\hand{289}}
